% TOP_PROBLEM: {"id": 30006897, "title": "Lieb-Thirring Conjecture on the Sharp Constant for Riesz Means of Order One in Three Dimensions", "category_id": 8, "category": {"id": 8, "name": "analysis", "display_name": "Analysis", "description": "Limits, continuity, calculus, and function theory.", "slug": "analysis", "order_index": 8, "created_at": "2026-07-31T15:26:25.671Z"}, "status": "open"}
% FIRST_POSTED: 2026-09-25
% ENTRY_KIND: statement_only
% !TeX program = lualatex
% Definition and sources only; this file is not an LLM solution attempt.
\documentclass[11pt]{article}
\usepackage[margin=25mm]{geometry}
\usepackage{fontspec}
\setmainfont{Latin Modern Roman}
\usepackage{amsmath,amssymb,mathtools}
\usepackage{unicode-math}
\setmathfont{Latin Modern Math}
\usepackage{xurl,hyperref}
\hypersetup{colorlinks=true,urlcolor=blue}
\setcounter{secnumdepth}{0}
\setlength{\parindent}{0pt}
\setlength{\parskip}{5pt}
\setlength{\emergencystretch}{3em}
\begin{document}
\section*{\texorpdfstring{Lieb-Thirring Conjecture on the Sharp Constant for Riesz Means of Order One in Three Dimensions}{Lieb-Thirring Conjecture on the Sharp Constant for Riesz Means of Order One in Three Dimensions}}\label{30006897}
\subsection{Definitions and mathematical statement}
For a real potential $V$ on ${\mathbb{R}}^3$, let $V_- =\max\{-V,0\}$ and form the self-adjoint Schr\"odinger operator $H=-\Delta+V$ through its quadratic form. Its negative spectral part is $H_-=\max\{-H,0\}$, so $\operatorname{Tr}H_-$ is the sum of magnitudes of negative eigenvalues, with multiplicity. The assertion is
\[
 \operatorname{Tr}(-\Delta+V)_-\le
 \frac1{15\pi^2}\int_{{\mathbb{R}}^3}V_-^{5/2}{\,\mathrm{d}} x,
\]
for the admissible potential class of the Lieb--Thirring inequality, with this constant optimal. It suffices for the constant question to formulate the supremum over smooth compactly supported negative potentials and use the standard form approximation extension. The exponent is the order-one, three-dimensional case; other Riesz orders or dimensions have different sharp constants.
\par\smallskip\textit{Formulation and concept references:} \href{https://arxiv.org/abs/2007.09326}{[S1]}.

\subsection{Short English statement}
Is the semiclassical value $1/(15\pi^2)$ the exact best constant controlling total negative Schr\"odinger energy by the integral of the potential's negative part to power $5/2$?
\subsection{Sources}
\begin{itemize}
\item[S1]Rupert L. Frank, The Lieb-Thirring inequalities: Recent results and open problems, arXiv:2007.09326 (2020).
\url{https://arxiv.org/abs/2007.09326}.
\textit{Formulation source.}
\end{itemize}
\end{document}
